%============================================================
%  Bd. 23 - Endliche Halbgruppen
%============================================================

\documentclass[main.tex]{subfiles}

\ifSubfilesClassLoaded{
  \BandSetupStandalone{B23}
}{
}

\title{Bd. 23 - Endliche Halbgruppen}
\author{Martin Kunze}
\date{}
\setcounter{file}{23}

\begin{document}

\title{Bd. 23 - Endliche Halbgruppen}
\author{Martin Kunze}
\date{}
\hypersetup{pageanchor=false}
\maketitle
\hypersetup{pageanchor=true}
\setcounter{tocdepth}{3}
\tableofcontents
\setcounter{file}{23}
\renewcommand{\FormulaBandID}{23}

\chapter{Einführung}

Dieser Band untersucht Halbgruppen mit endlichem Träger und das Isomorphieproblem ihrer Potenzhalbgruppen.

\chapter{Endliche Halbgruppen}

\section{Definition}

\FormulaDefDeltaK[Begriff der endlichen Halbgruppe]{%
  \FiniteSemigroupAlg{A}{\star}%
}{FiniteSemigroupDef}{%
  \DeltaRow{Mengen}{A}
  \DeltaRow{\textbf{Axiome}}{}
  \DeltaPrem{Halbgruppe}{\SemigroupAlg{A}{\star}}
  \DeltaPrem{Endlichkeit}{\Finite{A}}
  \DeltaRow{\textbf{Neues Symbol}}{}
  \DeltaPrem{Endliche Halbgruppe}{\FiniteSemigroupAlg{A}{\star}}
}

\FormulaAxiomDeltaK[Zugriff auf die Halbgruppenstruktur]{%
  \FiniteSemigroupAlg{A}{\star}\vdash\SemigroupAlg{A}{\star}%
}{FiniteSemigroupBaseAxiom}{%
  \DeltaRow{Mengen}{A}
}

\FormulaAxiomDeltaK[Zugriff auf die Endlichkeit]{%
  \FiniteSemigroupAlg{A}{\star}\vdash\Finite{A}%
}{FiniteSemigroupFinitenessAxiom}{%
  \DeltaRow{Mengen}{A}
}

\section{Beispiele}

\FormulaThmDeltaK[Eine konstante Operation bildet auf einer Paarmenge eine
endliche Halbgruppe]{%
  \begin{aligned}[t]
    &\forall x,y\in\{p,q\}\,(x\star_0y=p)\vdash{}\\[-2pt]
    &\FiniteSemigroupAlg{\{p,q\}}{\star_0}
  \end{aligned}
}{ConstantFinitePairSemigroup}{%
  \DeltaRow{Mengen}{p\dsep q\dsep x\dsep y}
  \DeltaRow{Funktionen}{\star_0}
}
\begin{tabproof}
  \proofstep{1}{\forall x,y\in\{p,q\}\,(x\star_0y=p)}{\rA}
  \proofstep{}{p\in\{p,q\}}{\FormulaRefAuto{a\in\{a,b\}}}
  \proofstep{1}{\SemigroupWithZeroAlg{\{p,q\}}{\star_0}{p}}{%
    \FormulaRefAuto{ConstantSemigroupWithZero}[thm]{2,1}}
  \proofstep{1}{\SemigroupAlg{\{p,q\}}{\star_0}}{%
    \FormulaRefAuto{SemigroupWithZeroBaseAxiom}[ax]{3}}
  \proofstep{}{\Finite{\{p,q\}}}{\FormulaRefAuto{FinitePair}[thm]}
  \proofstep{1}{\FiniteSemigroupAlg{\{p,q\}}{\star_0}}{%
    \FormulaRefAuto{FiniteSemigroupDef}[def]{4,5}}
\end{tabproof}

Für $p\neq q$ besitzt der Träger dieses Beispiels genau zwei Elemente.

\section{Endlichkeit der Potenzhalbgruppe}

Der folgende Satz verbindet die zuvor eingeführte endliche Spezialisierung
mit dem Potenzhalbgruppenbeispiel.

\FormulaThmDeltaK[Potenzhalbgruppen endlicher Halbgruppen sind endlich]{%
  \FiniteSemigroupAlg{A}{\star}
  \vdash\FiniteSemigroupAlg{\PowNE{A}}{\star_{\mathcal P}}%
}{FinitePowerSemigroup}{%
  \DeltaRow{Mengen}{A}
  \DeltaRow{Funktionen}{\star\dsep\star_{\mathcal P}}
}
\begin{tabproof}
  \proofstep{1}{\FiniteSemigroupAlg{A}{\star}}{\rA}
  \proofstep{1}{\SemigroupAlg{A}{\star}}{%
    \FormulaRefAuto{FiniteSemigroupBaseAxiom}[ax]{1}}
  \proofstep{1}{\Finite{A}}{%
    \FormulaRefAuto{FiniteSemigroupFinitenessAxiom}[ax]{1}}
  \proofstep{1}{\SemigroupAlg{\PowNE{A}}{\star_{\mathcal P}}}{%
    \FormulaRefAuto{NonemptyPowerSemigroup}[thm]{2}}
  \proofstep{}{\PowNE{A}\subseteq\powerset(A)}{%
    \FormulaRefAuto{A\setminus B\subseteq A}[thm]}
  \proofstep{1}{\Finite{\powerset(A)}}{%
    \FormulaRefAuto{FinitePowerSet}[thm]{3}}
  \proofstep{1}{\Finite{\PowNE{A}}}{%
    \FormulaRefAuto{FiniteSubset}[thm]{5,6}}
  \proofstep{1}{\FiniteSemigroupAlg{\PowNE{A}}{\star_{\mathcal P}}}{%
    \FormulaRefAuto{FiniteSemigroupDef}[def]{4,7}}
\end{tabproof}

\section{Die Potenzhalbgruppenvermutung}

Die globale Isomorphie und die allgemeine Richtung des Isomorphieproblems
wurden in Band~15 eingeführt. Das zugehörige Isomorphieproblem geht auf
Arbeiten von Tamura und Shafer zurück.

\begin{samepage}
\begin{center}
\fbox{\begin{minipage}{0.9\linewidth}
\textbf{Potenzhalbgruppenvermutung
  (Isomorphieproblem für endliche Potenzhalbgruppen).}
Für endliche Halbgruppen \((A,\star)\) und \((B,\diamond)\) gilt
\[
\begin{aligned}
&\exists F\,
  \SemigroupIso{F}
    {\PowNE{A},\star_{\mathcal P}}
    {\PowNE{B},\diamond_{\mathcal P}}
\quad\Longrightarrow{}\\[-2pt]
&\hspace{35mm}
  \exists f\,\SemigroupIso{f}{A,\star}{B,\diamond}.
\end{aligned}
\]
Formal stehen dabei
\(\FiniteSemigroupAlg{A}{\star}\) und
\(\FiniteSemigroupAlg{B}{\diamond}\) als Voraussetzungen fest.
\end{minipage}}
\end{center}
\end{samepage}

Die in Band~15 bewiesene umgekehrte Richtung gilt sogar ohne Endlichkeit.  Nur
die eingerahmte Richtung wird bewusst weder als Axiom noch als Theorem
registriert.

Ohne Endlichkeitsvoraussetzung ist die entsprechende Aussage falsch.  Das
vollständig bewiesene unendliche Gegenbeispiel steht bereits in
\FormulaRefAuto{MogiljanskajaPairCounterexample}[thm].  Es berührt die
hier formulierte Vermutung für endliche Halbgruppen nicht.

\subsection{Bemerkung: Bekannte Resultate und Status}

\begin{remark}
Die Vermutung ist mit Stand vom 1.~August~2026 für beliebige endliche
Halbgruppen offen; siehe
\href{https://doi.org/10.1007/978-3-031-75326-8_20}%
  {Tringali, \emph{On the Isomorphism Problem for Power Semigroups}, 2025}
und
\href{https://doi.org/10.48550/arXiv.2606.01917}%
  {Liu--Tringali, \emph{Power Semigroups and Two Rigidity Theorems for
  Groups}, 2026}.
Tamuras früherer Beitrag
\href{https://doi.org/10.1007/978-94-009-3839-7_22}%
  {\emph{On the Recent Results in the Study of Power Semigroups}, 1987}
kündigt die positive Aussage ausdrücklich ohne vollständigen Beweis an und
wird hier daher nicht als Lösung verwendet.

Wichtige bekannte Teilresultate sind:
\begin{itemize}
  \item Ohne die Endlichkeitsvoraussetzung ist die allgemeine Aussage nach
    \FormulaRefAuto{MogiljanskajaPairCounterexample}[thm] falsch; die
    verwendete Konstruktion geht zurück auf
    \href{https://doi.org/10.1007/BF02389140}%
      {Mogiljanskaja, \emph{Non-isomorphic Semigroups with Isomorphic
      Semigroups of Subsets}, 1973}.
  \item Sind beide Halbgruppen kürzbar und ist mindestens eine von ihnen
    kommutativ, so gilt das Isomorphieprinzip
    \href{https://doi.org/10.1007/978-3-031-75326-8_20}%
      {(Tringali, 2025)}.
  \item Die Klasse der idempotenten Halbgruppen ist global bestimmt
    \href{https://doi.org/10.1080/00927872.2017.1319474}%
      {(Yu--Zhao--Gan, 2018)}.
  \item Die Klasse der vollständig regulären Halbgruppen ist global bestimmt
    \href{https://doi.org/10.1080/00927872.2026.2659846}%
      {(Yu--Zhao, 2026)}.
  \item Ist eine der beiden Strukturen eine Gruppe, so braucht die andere
    nicht einmal vorab als Gruppe vorausgesetzt zu werden
    \href{https://doi.org/10.48550/arXiv.2606.01917}%
      {(Liu--Tringali, 2026)}.
\end{itemize}
Diese Spezialfälle entscheiden die Vermutung für endliche Halbgruppen im
Allgemeinen nicht.
\end{remark}

\section{Morphismen}

\FormulaDefDeltaK[Homomorphismus endlicher Halbgruppen]{%
  \FiniteSemigroupHom{f}{A,\star}{B,\diamond}%
}{FiniteSemigroupHomDef}{%
  \DeltaRow{Mengen}{A\dsep B}
  \DeltaRow{Funktionen}{f\dsep\star\dsep\diamond}
  \DeltaRow{\textbf{Axiome}}{}
  \DeltaPrem{Quellstruktur}{\FiniteSemigroupAlg{A}{\star}}
  \DeltaPrem{Zielstruktur}{\FiniteSemigroupAlg{B}{\diamond}}
  \DeltaPrem{Halbgruppenhomomorphismus}{%
    \SemigroupHom{f}{A,\star}{B,\diamond}}
  \DeltaRow{\textbf{Neues Symbol}}{}
  \DeltaPrem{Homomorphismus}{%
    \FiniteSemigroupHom{f}{A,\star}{B,\diamond}}
}

\FormulaAxiomDeltaK[Quellstruktur]{%
  \FiniteSemigroupHom{f}{A,\star}{B,\diamond}
  \vdash\FiniteSemigroupAlg{A}{\star}%
}{FiniteSemigroupHomSourceAxiom}{%
  \DeltaRow{Mengen}{A\dsep B}
  \DeltaRow{Funktionen}{f\dsep\star\dsep\diamond}
}

\FormulaAxiomDeltaK[Zielstruktur]{%
  \FiniteSemigroupHom{f}{A,\star}{B,\diamond}
  \vdash\FiniteSemigroupAlg{B}{\diamond}%
}{FiniteSemigroupHomTargetAxiom}{%
  \DeltaRow{Mengen}{A\dsep B}
  \DeltaRow{Funktionen}{f\dsep\star\dsep\diamond}
}

\FormulaAxiomDeltaK[Zugriff auf den Halbgruppenhomomorphismus]{%
  \FiniteSemigroupHom{f}{A,\star}{B,\diamond}
  \vdash\SemigroupHom{f}{A,\star}{B,\diamond}%
}{FiniteSemigroupHomBaseHomAxiom}{%
  \DeltaRow{Mengen}{A\dsep B}
  \DeltaRow{Funktionen}{f\dsep\star\dsep\diamond}
}

\FormulaDefDeltaK[Isomorphismus endlicher Halbgruppen]{%
  \FiniteSemigroupIso{f}{A,\star}{B,\diamond}%
}{FiniteSemigroupIsoDef}{%
  \DeltaRow{Mengen}{A\dsep B}
  \DeltaRow{Funktionen}{f\dsep\star\dsep\diamond}
  \DeltaRow{\textbf{Axiome}}{}
  \DeltaPrem{Homomorphismus}{%
    \FiniteSemigroupHom{f}{A,\star}{B,\diamond}}
  \DeltaPrem{Bijektivität}{f\colon A\bij B}
  \DeltaRow{\textbf{Neues Symbol}}{}
  \DeltaPrem{Isomorphismus}{%
    \FiniteSemigroupIso{f}{A,\star}{B,\diamond}}
}

\FormulaAxiomDeltaK[Homomorphismuszugriff]{%
  \FiniteSemigroupIso{f}{A,\star}{B,\diamond}
  \vdash\FiniteSemigroupHom{f}{A,\star}{B,\diamond}%
}{FiniteSemigroupIsoHomAxiom}{%
  \DeltaRow{Mengen}{A\dsep B}
  \DeltaRow{Funktionen}{f\dsep\star\dsep\diamond}
}

\FormulaAxiomDeltaK[Bijektivitätszugriff]{%
  \FiniteSemigroupIso{f}{A,\star}{B,\diamond}
  \vdash f\colon A\bij B%
}{FiniteSemigroupIsoBijectiveAxiom}{%
  \DeltaRow{Mengen}{A\dsep B}
  \DeltaRow{Funktionen}{f\dsep\star\dsep\diamond}
}

\section{Identifikation durch Multiplikationstafeln}

\subsection{Nummerierte Multiplikationstafeln}

Ist \(\FiniteSemigroupAlg{A}{\star}\), so liefert die
Bijektionscharakterisierung der Endlichkeit eine natürliche Zahl \(n\) und
eine Nummerierung
\[
  e\colon\NatSeg{n}\bij A.
\]
Dabei beginnt die Nummerierung mit dem Index \(0\).  Die Verknüpfung auf
\(A\) lässt sich mit \(e\) auf den Peano-Anfangsabschnitt
\(\NatSeg{n}\) zurücktransportieren.

\FormulaDefDeltaK[Nummerierte Multiplikationstafel einer endlichen
Halbgruppe]{%
  m_{\star,e}(i,j)
  \coloneqq e^{-1}\!\bigl(e(i)\star e(j)\bigr)%
}{FiniteSemigroupMultiplicationTableDef}{%
  \DeltaRow{Mengen}{A\dsep n\dsep i\dsep j}
  \DeltaRow{Funktionen}{e\dsep\star\dsep m_{\star,e}}
  \DeltaRow{\textbf{Axiome}}{}
  \DeltaPrem{Endliche Halbgruppe}{\FiniteSemigroupAlg{A}{\star}}
  \DeltaPrem{Länge der Nummerierung}{n\in\mathbb N}
  \DeltaPrem{Nummerierung}{e\colon\NatSeg{n}\bij A}
  \DeltaPrem{Indizes}{i,j\in\NatSeg{n}}
  \DeltaRow{\textbf{Neues Symbol}}{}
  \DeltaPrem{Tafeleintrag}{m_{\star,e}(i,j)}
}

Die Abgeschlossenheit von \(\star\) und die Bijektivität von \(e\)
stellen sicher, dass jeder Tafeleintrag wieder in \(\NatSeg{n}\) liegt.
Damit ist \(m_{\star,e}\) eine binäre Operation auf
\(\NatSeg{n}\).  Ihre Werte sind genau die Einträge der
Multiplikationstafel: Der erste Index bezeichnet die Zeile, der zweite die
Spalte.

\subsection{Das exakte Umbenennungskriterium}

Auf einem bereits nummerierten Träger ist eine bijektive Abbildung genau
dann ein Isomorphismus, wenn sie alle Tafeleinträge erhält.

\FormulaThmDeltaK[Umbenennungskriterium für endliche
Halbgruppentafeln]{%
  \begin{aligned}[t]
    &n\in\mathbb N\dsep
      \FiniteSemigroupAlg{\NatSeg{n}}{\star}\dsep
      \FiniteSemigroupAlg{\NatSeg{n}}{\diamond}\dsep
      \pi\colon\NatSeg{n}\bij\NatSeg{n}\vdash{}\\[-2pt]
    &\FiniteSemigroupIso{\pi}
      {\NatSeg{n},\star}{\NatSeg{n},\diamond}
      \leftrightarrow
      \forall i,j\in\NatSeg{n}\,
        \bigl(\pi(i\star j)=\pi(i)\diamond\pi(j)\bigr)
  \end{aligned}%
}{FiniteSemigroupTableRelabelCriterion}{%
  \DeltaRow{Mengen}{n\dsep i\dsep j}
  \DeltaRow{Funktionen}{\pi\dsep\star\dsep\diamond}
}
\begin{tabproofsplitwide}
  \proofpartwide{Hinrichtung}
    \proofstepwidestar[1]{n\in\mathbb N}{\rA}
    \proofstepwidestar[2]{%
      \FiniteSemigroupAlg{\NatSeg{n}}{\star}}{\rA}
    \proofstepwidestar[3]{%
      \FiniteSemigroupAlg{\NatSeg{n}}{\diamond}}{\rA}
    \proofstepwidestar[4]{%
      \pi\colon\NatSeg{n}\bij\NatSeg{n}}{\rA}
    \proofstepwidestar[5]{%
      \FiniteSemigroupIso{\pi}
        {\NatSeg{n},\star}{\NatSeg{n},\diamond}}{\rA}
    \proofstepwidestar[5]{%
      \FiniteSemigroupHom{\pi}
        {\NatSeg{n},\star}{\NatSeg{n},\diamond}}{%
      \FormulaRefAuto{FiniteSemigroupIsoHomAxiom}[ax]{5}}
    \proofstepwidestar[5]{%
      \SemigroupHom{\pi}
        {\NatSeg{n},\star}{\NatSeg{n},\diamond}}{%
      \FormulaRefAuto{FiniteSemigroupHomBaseHomAxiom}[ax]{6}}
    \proofstepwidestar[8]{i\in\NatSeg{n}}{\rA}
    \proofstepwidestar[9]{j\in\NatSeg{n}}{\rA}
    \proofstepwidestar[5,8,9]{%
      \pi(i\star j)=\pi(i)\diamond\pi(j)}{%
      \FormulaRefAuto{SemigroupHomOperationAxiom}[ax]{7,8,9}}
    \proofstepwidestar[5,8]{%
      \forall j\in\NatSeg{n}\,
        \bigl(\pi(i\star j)=\pi(i)\diamond\pi(j)\bigr)}{%
      \rUI{\rRI{9,10}}}
    \proofstepwidestar[5]{%
      \forall i,j\in\NatSeg{n}\,
        \bigl(\pi(i\star j)=\pi(i)\diamond\pi(j)\bigr)}{%
      \rUI{\rRI{8,11}}}
  \closeproofpartwide

  \proofpartwide{Rückrichtung}
    \proofstepwidestar[1]{n\in\mathbb N}{\rA}
    \proofstepwidestar[2]{%
      \FiniteSemigroupAlg{\NatSeg{n}}{\star}}{\rA}
    \proofstepwidestar[3]{%
      \FiniteSemigroupAlg{\NatSeg{n}}{\diamond}}{\rA}
    \proofstepwidestar[4]{%
      \pi\colon\NatSeg{n}\bij\NatSeg{n}}{\rA}
    \proofstepwidestar[5]{%
      \forall i,j\in\NatSeg{n}\,
        \bigl(\pi(i\star j)=\pi(i)\diamond\pi(j)\bigr)}{\rA}
    \proofstepwidestar[6]{i\in\NatSeg{n}}{\rA}
    \proofstepwidestar[7]{j\in\NatSeg{n}}{\rA}
    \proofstepwidestar[5,6]{%
      \forall j\in\NatSeg{n}\,
        \bigl(\pi(i\star j)=\pi(i)\diamond\pi(j)\bigr)}{%
      \rRE{\rUE{5},6}}
    \proofstepwidestar[5,6,7]{%
      \pi(i\star j)=\pi(i)\diamond\pi(j)}{%
      \rRE{\rUE{8},7}}
    \proofstepwidestar[2]{\SemigroupAlg{\NatSeg{n}}{\star}}{%
      \FormulaRefAuto{FiniteSemigroupBaseAxiom}[ax]{2}}
    \proofstepwidestar[3]{\SemigroupAlg{\NatSeg{n}}{\diamond}}{%
      \FormulaRefAuto{FiniteSemigroupBaseAxiom}[ax]{3}}
    \proofstepwidestar[4]{\pi\colon\NatSeg{n}\to\NatSeg{n}}{%
      \FormulaRefAuto{Bijektive Funktion}[def]{4}}
    \proofstepwidestar[2,3,4,5]{%
      \SemigroupHom{\pi}
        {\NatSeg{n},\star}{\NatSeg{n},\diamond}}{%
      \FormulaRefAuto{SemigroupHomDef}[def]{10,11,12,9}}
    \proofstepwidestar[2,3,4,5]{%
      \FiniteSemigroupHom{\pi}
        {\NatSeg{n},\star}{\NatSeg{n},\diamond}}{%
      \FormulaRefAuto{FiniteSemigroupHomDef}[def]{2,3,13}}
    \proofstepwidestar[2,3,4,5]{%
      \FiniteSemigroupIso{\pi}
        {\NatSeg{n},\star}{\NatSeg{n},\diamond}}{%
      \FormulaRefAuto{FiniteSemigroupIsoDef}[def]{14,4}}
  \closeproofpartwide
\end{tabproofsplitwide}

Nun seien \(e\colon\NatSeg{n}\bij A\) und
\(d\colon\NatSeg{n}\bij B\) Nummerierungen zweier endlicher
Halbgruppen gleicher Trägergröße.  Auf ihre nummerierten Tafeln
angewandt ergibt der Satz das vollständige Identifikationskriterium
\[
\begin{aligned}
&\exists f\,
  \FiniteSemigroupIso{f}{A,\star}{B,\diamond}
\quad\Longleftrightarrow{}\\[-2pt]
&\exists\pi\Bigl(
  \pi\colon\NatSeg{n}\bij\NatSeg{n}\ \land
  \forall i,j\in\NatSeg{n}\,
  m_{\diamond,d}\bigl(\pi(i),\pi(j)\bigr)
    =\pi\bigl(m_{\star,e}(i,j)\bigr)
  \Bigr).
\end{aligned}
\]
Es muss also dieselbe Umbenennung auf die Zeilen, die Spalten und die
Tafeleinträge angewandt werden.  In der Hinrichtung ist
\(\pi=d^{-1}\circ f\circ e\); in der Rückrichtung gewinnt man den
Isomorphismus als \(f=d\circ\pi\circ e^{-1}\).  Dies beweist zugleich,
dass das Kriterium notwendig und hinreichend ist.

\subsection{Ein kanonischer Tafelcode}

Das Umbenennungskriterium liefert einen exakten, endlichen
Identifikationsalgorithmus.  Für eine Tafel \(m\) auf \(\NatSeg{n}\)
und eine Selbstbijektion
\(\pi\colon\NatSeg{n}\bij\NatSeg{n}\) sei
\[
  m^{\pi}(i,j)
  \coloneqq
  \pi\!\left(m\bigl(\pi^{-1}(i),\pi^{-1}(j)\bigr)\right).
\]
Man liest die Werte von \(m^{\pi}\) zeilenweise als ein Wort
\(w(m^{\pi})\) der Länge \(n^2\); für \(n=0\) ist dies das leere
Wort.  Mit der natürlichen lexikographischen Ordnung dieser Wörter setze
man, für eine beliebig gewählte Nummerierung \(e\),
\[
  \operatorname{can}(A,\star)
  \coloneqq
  \left(
    n,
    \min_{\mathrm{lex}}
    \left\{
      w\!\left(m_{\star,e}^{\pi}\right)
      \;\middle|\;
      \pi\colon\NatSeg{n}\bij\NatSeg{n}
    \right\}
  \right).
\]
Es gibt nur endlich viele Selbstbijektionen, daher existiert das Minimum.
Eine andere Ausgangsnummerierung verändert nur die Reihenfolge der Tafeln
in der Menge, nicht ihr Minimum.  Das vorangestellte \(n\) unterscheidet
auch leere Tafelwörter verschiedener formaler Herkunft.  Aus dem
Umbenennungskriterium folgt die vollständige Charakterisierung
\[
  (A,\star)\cong(B,\diamond)
  \quad\Longleftrightarrow\quad
  \operatorname{can}(A,\star)
    =\operatorname{can}(B,\diamond).
\]

Für einen Träger mit \(n\) Elementen sind höchstens \(n!\)
Umbenennungen und je Umbenennung \(n^2\) Tafeleinträge zu prüfen.  Bei
zwei beziehungsweise drei Elementen sind dies nur \(2\) beziehungsweise
\(6\) Umbenennungen.  Als schnelle Vorfilter kann man etwa die Anzahl der
idempotenten Elemente, Kommutativität, die Existenz eines Nullelements,
Kürzbarkeit und die in Band~15 eingeführten Teilbarkeitsrelationen
vergleichen.  Verschiedene Werte schließen einen Isomorphismus aus;
übereinstimmende Werte ersetzen den kanonischen Tafelcode im Allgemeinen
nicht.

\begin{remark}[Grenze für die Potenzhalbgruppenvermutung]
Das Tafelverfahren identifiziert endliche Halbgruppen, sobald ihre
Multiplikationstafeln vorliegen.  Aus einem abstrakten Isomorphismus ihrer
Potenzhalbgruppen lässt sich jedoch nicht ohne Weiteres eine Umbenennung der
Grundträger ablesen: Einermengen sind rein multiplikativ nicht im
Allgemeinen ausgezeichnet.

Schon für eine Nullhalbgruppe mit mindestens zwei Elementen gilt für alle
nichtleeren Teilmengen \(X,Y\) stets
\(X\star_{\mathcal P}Y=\{0\}\).  Daher ist jede Permutation von
\(\PowNE{A}\setminus\{\{0\}\}\), welche \(\{0\}\) fest lässt, ein
Automorphismus der Potenzhalbgruppe; eine solche Permutation kann etwa eine
Einermenge mit einer zweielementigen Teilmenge vertauschen.  Genau deshalb
ist der kanonische Tafelcode ein vollständiger Isomorphietest für gegebene
Tafeln, aber noch kein allgemeiner Beweis der Potenzhalbgruppenvermutung.
\end{remark}

\end{document}
